\documentclass{article}
\usepackage[utf8]{inputenc}
\usepackage[top=4.5cm, bottom=2cm, left=2cm, right=2cm, headheight=4.5cm]{geometry}
\usepackage{background}
\usepackage{fancyhdr}
\usepackage{graphicx}
\usepackage{array}
\usepackage{calc}
\usepackage[hidelinks]{hyperref}
\usepackage{enumerate}
\usepackage{enumitem}
\usepackage{pdfpages}
\usepackage{amsmath}
\usepackage{amssymb}
\usepackage{cleveref}
\usepackage{svg}
\usepackage{stackrel}
\usepackage{algorithm}
\usepackage[noend]{algpseudocode}
\usepackage{listings}
\newcommand{\acados}{\texttt{acados}}
\newcommand{\leapc}{\texttt{leap-c}}

% \usepackage[framed,numbered,autolinebreaks,useliterate]{mcode}
\usepackage[short]{optidef}
\newcommand{\norm}[1]{\left\lVert#1\right\rVert}

\backgroundsetup{
  scale=1,
  opacity=0.7, % Adjust the opacity as needed
%   position={2cm, 0cm},
  placement=top,
  contents={\includegraphics[trim=0 3.0cm 0 0, clip, width=\textwidth]{img/title.png}}  % Adjust the width as needed}
}

\pagestyle{fancy}
\fancyhf{}  % Clear the existing header and footer settings
\chead{
	\begin{center}
		{\Large
		\textsc{Model Predictive Control and Reinforcement Learning}\\
		\vspace*{0.3cm}
		Exercise Sheet 2: Introduction to \acados{} and \leapc{} in \texttt{Python}}
	\end{center}
	\hfill{\small Joschka Boedecker, Moritz Diehl, Leonard Fichtner, Jonathan Frey, Andrea Ghezzi, Jasper Hoffmann}
}
\rfoot{\thepage}

%opening
\begin{document}



\section*{Exercise 1: Getting started with \acados{}}
In this exercise, you will implement a simple MPC controller for the Cart Pole environment
(aka inverted pendulum on a cart), see Figure \ref{cartpole}.
The goal is to implement a controller that balances the pole in the closed-loop simulation provided in \verb*|example_closed_loop.py|.

\begin{figure}[ht!]
\center
\includegraphics[width=0.7\textwidth]{img/cartpole.png}
\caption{An illustration of Cart Pole while swinging up.}
\label{cartpole}
\end{figure}


\begin{enumerate}[label=1.\arabic*]
	\item Familiarize yourself with the \href{https://github.com/leap-c/leap-c/blob/main/leap_c/examples/cartpole/env.py#L24}{\texttt{CartPoleEnv}} environment, e.g., by reading the documentation of the class.
	\item Complete the controller formulation in the function \verb*|setup_ocp| provided in \verb*|setup_ocp.py|.
	For a given initial state $s$, this function formulates a simple MPC controller described by the problem formulation
	\begin{align*}
    \min_{\substack{x_0, u_0, x_1,...,u_{N-1},x_N}} & \frac{1}{N}\sum_{k=0}^{N-1} l(x_k, u_k) + \overline{V}(x_N)\\
                                                   \text{s.t. } & x_0 = s\\
                                                    & x_{k+1} = f^{\mathrm{dyn}}(x_k, u_k),\ k = 0,...,N-1\\
                                                    & \underline{u} \leq u_k \leq \overline{u},\ k = 0,...,N-1,
	\end{align*}
	where $f^{\mathrm{dyn}}(x_k, u_k)$ are the dynamics of the Cart Pole system, given by one step of an RK4 integrator, $\underline{u}$ and
	$\overline{u}$ are the actuator limits, and the stage costs and terminal cost are given by
	\begin{equation*} l(x_k, u_k)= \frac{1}{2}||\begin{pmatrix}x_k\\u_k\end{pmatrix}-\begin{pmatrix}x^{\mathrm{ref}}\\u^\mathrm{ref}\end{pmatrix}||_{W}^2\end{equation*}
	and
	\begin{equation*} \overline{V}(x_N)= \frac{1}{2}||x_N-x_{\mathrm{e}}^\mathrm{ref}||_{W_{\mathrm{e}}}^2 .\end{equation*}

	Set the values of $x^{\mathrm{ref}}, u^{\mathrm{ref}}, x^{\mathrm{ref}}_{\mathrm{e}}, W, W_\mathrm{e}$ in a way that the controller successfully balances the pole.
	This will be indicated by messages printed by \verb*|example_closed_loop.py|.
	This will also generate plots and a rendered video that might be helpful.
\end{enumerate}


\section*{Exercise 2: Parametric sensitivities}
In this exercise, we show how to extend the \acados{} formulation from the previous exercise step-by-step,
to arrive at a \leapc{} controller that can be used to train parameters with one of the
\leapc{} algorithms and \texttt{PyTorch} in general. As an example, we will
use this formulation to solve MPC problems for varying values for the reference of the angle $\theta$
and compute their sensitivities in \verb*|example_diffMPC.py|.

\begin{enumerate}[label=2.\arabic*]
	\item Make yourself familiar with the
	\verb*|AcadosParameterManager| class.
	\item The goal now is to make a modified version of \verb*|setup_ocp|,
	in which the parameter manager class is used to introduce new parameters into the
	acados ocp formulation.\\
	At the bottom of \verb*|setup_ocp_with_manager.py| you will find a method
	\verb*|create_cartpole_params| which returns a list of \verb*|AcadosParameters| to
	be passed to the \verb*|AcadosParameterManager| instance at creation.
	As an example, it already contains instances meant for the weight matrix.
	Extend this list to also contain 5 more \texttt{AcadosParameter}s, one for each field in the
	reference part of the cost.
	Make only the reference of the angle~$\theta$ learnable.
	\item Now complete the function \verb*|setup_ocp_with_manager|
	in \verb*|setup_ocp_with_manager.py|.
	\item Familiarize yourself with \verb*|AcadosDiffMpcTorch| and \\
	\verb*|AcadosController|.
	Create an \verb*|AcadosController| in \verb*|example_diffMPC.py|.
	Make sure to set \verb*|n_batch_max| of the \verb*|AcadosDiffMpcTorch| high enough
	such that the batch of parameters \verb*|p_values| can be passed.
	\item Run \verb*|example_diffMPC.py| to solve a batch
	of optimization problems with varying reference of the angle and compute their sensitivities.
	Think about what these plots say in the context of learning approaches aiming to leverage the sensitivities.
	\item (Bonus) Feel free to play around with values of the parameters, ocp formulations,
	calculating other sensitivities, or making parameters differ over the horizon 
	(field \verb*|end_stages| for the \verb|AcadosParameter|).
	Note that while playing around, it will likely happen, that the solver will not converge for some of the parameters you will throw at him.
	Therefore, you might want to adjust your ocp formulation (e.g., adjust ``solver\_options'', adjust fixed values, slack constraints),
	adjust the sampled parameters (e.g., their range), 
	or just ignore it and filter the non-converged samples for whatever you want to do (e.g., plotting).

\end{enumerate}

\subsection*{Useful Resources}
\begin{itemize}
	\item \href{https://github.com/acados/acados/blob/master/docs/problem_formulation/problem_formulation_ocp_mex.pdf}{\acados{} problem formulation PDF}
	\item \href{https://github.com/leap-c/leap-c/blob/main/leap_c/examples/cartpole/env.py#L24}{\texttt{CartPoleEnv}}
	\item \href{https://github.com/leap-c/leap-c/blob/main/leap_c/ocp/acados/parameters.py#L15}{\texttt{AcadosParameters}}
	\item \href{https://github.com/leap-c/leap-c/blob/main/leap_c/ocp/acados/parameters.py#L70}{\texttt{AcadosParameterManager}}
	\item \href{https://github.com/leap-c/leap-c/blob/main/leap_c/ocp/acados/torch.py#L18}{\texttt{AcadosDiffMpcTorch}}
	\item \href{https://github.com/leap-c/leap-c/blob/main/leap_c/ocp/acados/diff_mpc.py#L30}{\texttt{AcadosDiffMpcCtx}}
	\item \href{https://github.com/leap-c/leap-c/blob/main/leap_c/ocp/acados/controller.py#L13}{\texttt{AcadosController}}
\end{itemize}


%\section*{Bonus exercises}

\end{document}

